% A one-file article: everything a LaTeX document can need, carried inside it.
%
% This is the shape `05-SPEC-latex.md` was written against, before
% `08-SPEC-directories.md` made a document a directory. `filecontents*` still
% works and is what this file leans on, so the corpus keeps a case for it.
%
% It also carries, on purpose, the four things step 4's anchoring has to
% survive: a hyphenated line end, a page break inside a sentence, a footnote,
% and a ligature. Do not "tidy" the paragraph lengths below; they are the
% fixture.
\documentclass[11pt]{article}

\begin{filecontents*}[overwrite]{article.bib}
@article{knuth1984,
  author  = {Knuth, Donald E.},
  title   = {Literate Programming},
  journal = {The Computer Journal},
  volume  = {27},
  number  = {2},
  pages   = {97--111},
  year    = {1984},
}
@book{lamport1994,
  author    = {Lamport, Leslie},
  title     = {LaTeX: A Document Preparation System},
  publisher = {Addison-Wesley},
  year      = {1994},
}
\end{filecontents*}

\usepackage[T1]{fontenc}
\usepackage[utf8]{inputenc}
\usepackage{amsmath}
\usepackage{graphicx}

\title{A One-File Article}
\author{The LibrePaper Corpus}
\date{}

\begin{document}
\maketitle

\begin{abstract}
  One article in one file, with its bibliography riding along inside it, so
  that a compiler which has no files beside the document still has a
  bibliography to typeset.
\end{abstract}

\section{Introduction}

Literate programming asks that a program be written for a reader rather than
for a compiler~\cite{knuth1984}, and the same argument was later made about
the document preparation system this file is written in~\cite{lamport1994}.
The claim is not that either is easy but that the difficulty is
incompressible: it moves, it does not vanish.\footnote{A footnote, so the
anchoring has a text node that is nowhere near its mark on the page.}

The word at the end of this line is deliberately long and will be
hyphenated by \TeX{} rather than moved whole to the next line, because
incomprehensibility is not a word that fits in a column. That hyphen is a
character pdf.js will extract, and the anchor that spans it has to survive it.

The word ``efficient'' carries an \texttt{fi} ligature, which some fonts hand
to pdf.js as one glyph and some as two.

\section{Mathematics}

The exponential integral has no elementary antiderivative:

\begin{equation}
  \label{eq:ei}
  \operatorname{Ei}(x) = -\int_{-x}^{\infty} \frac{e^{-t}}{t}\,\mathrm{d}t,
  \qquad x > 0 .
\end{equation}

Equation~\eqref{eq:ei} is referenced here so the second pass has something to
resolve, which is how the log comes to ask for one.

\section{A figure with no image}

A figure environment with no \verb|\includegraphics| in it, because a
compiler with no files beside the document must still be able to typeset the
float, its caption and its number.

\begin{figure}[h]
  \centering
  \fbox{\parbox{0.6\textwidth}{\centering\vspace{2em}(no image)\vspace{2em}}}
  \caption{A float with nothing floating in it.}
  \label{fig:empty}
\end{figure}

\section{The page break}

% Tuned so TeX really breaks the sentence below across the page boundary
% rather than merely starting it near one: at 12em it fitted whole on the
% page and step 4 was testing nothing. Three pages either way.
\vspace{26em}

This sentence begins on one page and, because of the vertical space above it
and the filler below, is very likely to be broken across the page boundary by
\TeX{}, which is precisely the case a text-quote anchor has to re-anchor
across, since pdf.js will hand the reader two runs of text with a gap between
them where the page turned and nothing in the extracted sequence to say that
the gap is not a space.

\vspace{6em}

\section{Conclusion}

Nothing is concluded. The document exists to be compiled by three engines and
to produce three logs that agree about how many pages it has.

\bibliographystyle{plain}
\bibliography{article}

\end{document}
